% Options for packages loaded elsewhere
\PassOptionsToPackage{unicode}{hyperref}
\PassOptionsToPackage{hyphens}{url}
\PassOptionsToPackage{dvipsnames,svgnames,x11names}{xcolor}
%
\documentclass[
  11pt,
  a4paper,
]{article}

\usepackage{amsmath,amssymb}
\usepackage{setspace}
\usepackage{iftex}
\ifPDFTeX
  \usepackage[T1]{fontenc}
  \usepackage[utf8]{inputenc}
  \usepackage{textcomp} % provide euro and other symbols
\else % if luatex or xetex
  \usepackage{unicode-math}
  \defaultfontfeatures{Scale=MatchLowercase}
  \defaultfontfeatures[\rmfamily]{Ligatures=TeX,Scale=1}
\fi
\usepackage{lmodern}
\ifPDFTeX\else  
    % xetex/luatex font selection
  \setmainfont[]{Palatino Linotype}
\fi
% Use upquote if available, for straight quotes in verbatim environments
\IfFileExists{upquote.sty}{\usepackage{upquote}}{}
\IfFileExists{microtype.sty}{% use microtype if available
  \usepackage[]{microtype}
  \UseMicrotypeSet[protrusion]{basicmath} % disable protrusion for tt fonts
}{}
\makeatletter
\@ifundefined{KOMAClassName}{% if non-KOMA class
  \IfFileExists{parskip.sty}{%
    \usepackage{parskip}
  }{% else
    \setlength{\parindent}{0pt}
    \setlength{\parskip}{6pt plus 2pt minus 1pt}}
}{% if KOMA class
  \KOMAoptions{parskip=half}}
\makeatother
\usepackage{xcolor}
\usepackage[top=2.4cm,bottom=2.4cm,left=2.5cm,right=2.5cm]{geometry}
\setlength{\emergencystretch}{3em} % prevent overfull lines
\setcounter{secnumdepth}{-\maxdimen} % remove section numbering


\providecommand{\tightlist}{%
  \setlength{\itemsep}{0pt}\setlength{\parskip}{0pt}}\usepackage{longtable,booktabs,array}
\usepackage{calc} % for calculating minipage widths
% Correct order of tables after \paragraph or \subparagraph
\usepackage{etoolbox}
\makeatletter
\patchcmd\longtable{\par}{\if@noskipsec\mbox{}\fi\par}{}{}
\makeatother
% Allow footnotes in longtable head/foot
\IfFileExists{footnotehyper.sty}{\usepackage{footnotehyper}}{\usepackage{footnote}}
\makesavenoteenv{longtable}
\usepackage{graphicx}
\makeatletter
\def\maxwidth{\ifdim\Gin@nat@width>\linewidth\linewidth\else\Gin@nat@width\fi}
\def\maxheight{\ifdim\Gin@nat@height>\textheight\textheight\else\Gin@nat@height\fi}
\makeatother
% Scale images if necessary, so that they will not overflow the page
% margins by default, and it is still possible to overwrite the defaults
% using explicit options in \includegraphics[width, height, ...]{}
\setkeys{Gin}{width=\maxwidth,height=\maxheight,keepaspectratio}
% Set default figure placement to htbp
\makeatletter
\def\fps@figure{htbp}
\makeatother
% definitions for citeproc citations
\NewDocumentCommand\citeproctext{}{}
\NewDocumentCommand\citeproc{mm}{%
  \begingroup\def\citeproctext{#2}\cite{#1}\endgroup}
\makeatletter
 % allow citations to break across lines
 \let\@cite@ofmt\@firstofone
 % avoid brackets around text for \cite:
 \def\@biblabel#1{}
 \def\@cite#1#2{{#1\if@tempswa , #2\fi}}
\makeatother
\newlength{\cslhangindent}
\setlength{\cslhangindent}{1.5em}
\newlength{\csllabelwidth}
\setlength{\csllabelwidth}{3em}
\newenvironment{CSLReferences}[2] % #1 hanging-indent, #2 entry-spacing
 {\begin{list}{}{%
  \setlength{\itemindent}{0pt}
  \setlength{\leftmargin}{0pt}
  \setlength{\parsep}{0pt}
  % turn on hanging indent if param 1 is 1
  \ifodd #1
   \setlength{\leftmargin}{\cslhangindent}
   \setlength{\itemindent}{-1\cslhangindent}
  \fi
  % set entry spacing
  \setlength{\itemsep}{#2\baselineskip}}}
 {\end{list}}
\usepackage{calc}
\newcommand{\CSLBlock}[1]{\hfill\break\parbox[t]{\linewidth}{\strut\ignorespaces#1\strut}}
\newcommand{\CSLLeftMargin}[1]{\parbox[t]{\csllabelwidth}{\strut#1\strut}}
\newcommand{\CSLRightInline}[1]{\parbox[t]{\linewidth - \csllabelwidth}{\strut#1\strut}}
\newcommand{\CSLIndent}[1]{\hspace{\cslhangindent}#1}

\makeatletter
\@ifpackageloaded{caption}{}{\usepackage{caption}}
\AtBeginDocument{%
\ifdefined\contentsname
  \renewcommand*\contentsname{Table of contents}
\else
  \newcommand\contentsname{Table of contents}
\fi
\ifdefined\listfigurename
  \renewcommand*\listfigurename{List of Figures}
\else
  \newcommand\listfigurename{List of Figures}
\fi
\ifdefined\listtablename
  \renewcommand*\listtablename{List of Tables}
\else
  \newcommand\listtablename{List of Tables}
\fi
\ifdefined\figurename
  \renewcommand*\figurename{Figure}
\else
  \newcommand\figurename{Figure}
\fi
\ifdefined\tablename
  \renewcommand*\tablename{Table}
\else
  \newcommand\tablename{Table}
\fi
}
\@ifpackageloaded{float}{}{\usepackage{float}}
\floatstyle{ruled}
\@ifundefined{c@chapter}{\newfloat{codelisting}{h}{lop}}{\newfloat{codelisting}{h}{lop}[chapter]}
\floatname{codelisting}{Listing}
\newcommand*\listoflistings{\listof{codelisting}{List of Listings}}
\makeatother
\makeatletter
\makeatother
\makeatletter
\@ifpackageloaded{caption}{}{\usepackage{caption}}
\@ifpackageloaded{subcaption}{}{\usepackage{subcaption}}
\makeatother
\ifLuaTeX
  \usepackage{selnolig}  % disable illegal ligatures
\fi
\usepackage{bookmark}

\IfFileExists{xurl.sty}{\usepackage{xurl}}{} % add URL line breaks if available
\urlstyle{same} % disable monospaced font for URLs
\hypersetup{
  pdftitle={Psychological Empowerment and Turnover: A Review of the Literature},
  pdfauthor={Isaiah Espinoza},
  pdfkeywords={Psychological Empowerment, Turnover, Election
resiliency, Poll workers, Election administration},
  colorlinks=true,
  linkcolor={blue},
  filecolor={Maroon},
  citecolor={Blue},
  urlcolor={Blue},
  pdfcreator={LaTeX via pandoc}}

%% CAPTIONS
\usepackage{caption}
\DeclareCaptionStyle{italic}[justification=centering]
 {labelfont={bf},textfont={it},labelsep=colon}
\captionsetup[figure]{style=italic,format=hang,singlelinecheck=true}
\captionsetup[table]{style=italic,format=hang,singlelinecheck=true}

%% FONT
% \usepackage{bera}
% \usepackage[charter]{mathdesign}
% \usepackage[scale=0.9]{sourcecodepro}
% \usepackage[lf,t]{FiraSans}
\usepackage{fontawesome}

%% HEADERS AND FOOTERS
\usepackage{fancyhdr}
\pagestyle{fancy}
\rfoot{\Large\sffamily\raisebox{-0.1cm}{\textbf{\thepage}}}
\makeatletter
\lhead{\textsf{\expandafter{\@title}}}
\makeatother
\rhead{}
\cfoot{}
\setlength{\headheight}{15pt}
\renewcommand{\headrulewidth}{0.4pt}
\renewcommand{\footrulewidth}{0.4pt}
\fancypagestyle{plain}{%
\fancyhf{} % clear all header and footer fields
\fancyfoot[C]{\sffamily\thepage} % except the center
\renewcommand{\headrulewidth}{0pt}
\renewcommand{\footrulewidth}{0pt}}

%% MATHS
\usepackage{bm,amsmath}
\allowdisplaybreaks

%% GRAPHICS
\makeatletter
\def\fps@figure{htbp}
\makeatother
\setcounter{topnumber}{2}
\setcounter{bottomnumber}{2}
\setcounter{totalnumber}{4}
\renewcommand{\topfraction}{0.85}
\renewcommand{\bottomfraction}{0.85}
\renewcommand{\textfraction}{0.15}
\renewcommand{\floatpagefraction}{0.8}

%% SECTION TITLES
\usepackage[compact,sf,bf]{titlesec}
\titleformat*{\section}{\Large\sf\bfseries\color[rgb]{0.7,0,0}}
\titleformat*{\subsection}{\large\sf\bfseries\color[rgb]{0.7,0,0}}
\titleformat*{\subsubsection}{\sf\bfseries\color[rgb]{0.7,0,0}}
\titlespacing{\section}{0pt}{2ex}{.5ex}
\titlespacing{\subsection}{0pt}{1.5ex}{0ex}
\titlespacing{\subsubsection}{0pt}{.5ex}{0ex}


%% BIBLIOGRAPHY.

\makeatletter
\@ifpackageloaded{biblatex}{
\ExecuteBibliographyOptions{bibencoding=utf8,minnames=1,maxnames=3, maxbibnames=99,dashed=false,terseinits=true,giveninits=true,uniquename=false,uniquelist=false,doi=false, isbn=false,url=true,sortcites=false}
\DeclareFieldFormat{url}{\texttt{\url{#1}}}
\DeclareFieldFormat[article]{pages}{#1}
\DeclareFieldFormat[inproceedings]{pages}{\lowercase{pp.}#1}
\DeclareFieldFormat[incollection]{pages}{\lowercase{pp.}#1}
\DeclareFieldFormat[article]{volume}{\mkbibbold{#1}}
\DeclareFieldFormat[article]{number}{\mkbibparens{#1}}
\DeclareFieldFormat[article]{title}{\MakeCapital{#1}}
\DeclareFieldFormat[article]{url}{}
\DeclareFieldFormat[inproceedings]{title}{#1}
\DeclareFieldFormat{shorthandwidth}{#1}
\usepackage{xpatch}
\xpatchbibmacro{volume+number+eid}{\setunit*{\adddot}}{}{}{}
% Remove In: for an article.
\renewbibmacro{in:}{%
  \ifentrytype{article}{}{%
  \printtext{\bibstring{in}\intitlepunct}}}
\AtEveryBibitem{\clearfield{month}}
\AtEveryCitekey{\clearfield{month}}
\DeclareDelimFormat[cbx@textcite]{nameyeardelim}{\addspace}
\renewcommand*{\finalnamedelim}{\addspace\&\space}
}{}
\makeatother

%% PAGE BREAKING to avoid widows and orphans
\clubpenalty = 2000
\widowpenalty = 2000
\usepackage{microtype}% Placement of logos

\RequirePackage[absolute,overlay]{textpos}
\setlength{\TPHorizModule}{1cm}
\setlength{\TPVertModule}{1cm}
\def\placefig#1#2#3#4{\begin{textblock}{.1}(#1,#2)\rlap{\includegraphics[#3]{#4}}\end{textblock}}

% Title and date

\title{Psychological Empowerment and Turnover: A Review of the
Literature}
\date{April 22, 2024}

\def\Date{\number\day}
\def\Month{\ifcase\month\or
 January\or February\or March\or April\or May\or June\or
 July\or August\or September\or October\or November\or December\fi}
\def\Year{\number\year}

%%%% PAGE STYLE FOR FRONT PAGE OF REPORTS

\makeatletter
\def\organization#1{\gdef\@organization{#1}}
\def\telephone#1{\gdef\@telephone{#1}}
\def\email#1{\gdef\@email{#1}}
\makeatother
  \organization{}

  \def\name{Center for Democracy and Civic Engagement}
% 
  \telephone{(410) 903 6911}
% 
  \email{ssnovey@umd.edu}

\def\webaddress{\url{https://cdce.umd.edu/}}
\def\abn{12 377 614 012}
\def\extraspace{\vspace*{1.6cm}}
\makeatletter
\def\contactdetails{\faicon{phone} & \@telephone \\
                    \faicon{envelope} & \@email}
\makeatother

\usepackage[absolute,overlay]{textpos}
\setlength{\TPHorizModule}{1cm}
\setlength{\TPVertModule}{1cm}

%%%% FRONT PAGE OF REPORTS

% \def\reporttype{Report}

\long\def\front#1#2#3{
\newpage
% \begin{textblock}{7}(12.7,28.2)\hfill
% \includegraphics[height=4cm]{umd_logo2}~~~
% \end{textblock}
\begin{singlespacing}
\thispagestyle{empty}
\vspace*{-1.4cm}
\hspace*{-1.4cm}
\hbox to 16cm{
  \hbox to 6.5cm{\vbox to 14cm{\vbox to 25cm{
    \includegraphics[width=6.6cm]{CDCE}
    \vfill
    \includegraphics[width=5cm]{umd_logo2}
    \vspace{0.4cm}
    \par
    \parbox{6.3cm}{\raggedright
      \sf\color[rgb]{0.70,0.00,0.00}
      {\large\textbf{\name}}\par
      \vspace{.7cm}
      \tabcolsep=0.12cm\sf\small
      \begin{tabular}{@{}ll@{}}\contactdetails
      \end{tabular}
    }
  }\vss}\hss}
  \hspace*{0.2cm}
  \hbox to 1cm{\vbox to 14cm{\rule{1pt}{26.8cm}\vss}\hss\hfill}
  \hbox to 10cm{\vbox to 14cm{\vbox to 25cm{
      \vspace*{3cm}\sf\raggedright
      \parbox{10cm}{\sf\raggedright\baselineskip=1.2cm
         \fontsize{24.88}{30}\color[rgb]{0.70,0.00,0.00}\sf\textbf{#1}}
      \par
      \vfill
      \large
      \vbox{\parskip=0.8cm #2}\par
      \vspace*{2cm}\par
      % \reporttype\\[0.3cm]
      % \hbox{#3}%\\[2cm]\
      % \vspace*{1cm}
      %{\large\sf\textbf{\Date~\Month~\Year}}
   }\vss}
  }}
\end{singlespacing}
\newpage
}

\makeatletter
\def\maketitle{\front{\expandafter{\@title}}{\@author}{\@organization}}
\makeatother

% Authors

\author{\sf{\Large\textbf{Isaiah Espinoza} \\\large Research
Assistant\\[0.5cm]}}
%\lfoot{\sf Espinoza: April 22, 2024}
\begin{document}
\maketitle
\begin{abstract}
An increase in turnover among the election administration workforce has
been observed and become recently salient. Election administration work
(e.g., poll workers) is certainly difficult. The impetus of the
following review was to discover what lessons, supports, or specific
interventions could be derived from selected literature to potentially
help reduce turnover, turnover intention, and boost morale for election
workers in particular. A selection of academic literature on
psychological empowerment and turnover are reviewed for that purpose.
\end{abstract}

\setstretch{1.5}
\section{Introduction}\label{introduction}

Psychological empowerment has been studied extensively within
organizational psychology and management studies with the aim of
determining drivers for workplace performance and retention. Researchers
of the field were interested in understanding how workers adapt and
adjust to work environments, what motivates workers to go above and
beyond minimal workplace performance, and what might decrease the
likelihood of turnover. The reason for the interest in the individual
psychology of the worker was seemingly prompted by the limitations of
previous research up to that point. Prior to the shift toward the
psychological, research was limited to classical bureaucratic
explanations of how workplace performance and satisfaction was best
improved. That is, work was assumed to have only \emph{instrumental
value} to the workers. As such, performance, retention, and workplace
satisfaction was determined through different organizational controls in
combination with contingent rewards and punishments. As will be seen,
some researchers saw this tradition as limiting and were thus motivated
to determine and study those organizational structures and workplace
interventions that have the most empowering effects on the individual
workers themselves. This branch of research introduced a novel paradigm
exploring the significance of a worker's \emph{internalized commitment}
to the work itself in order to understand and help to foster worker
empowerment (\citeproc{ref-thomas1990}{Thomas and Velthouse 1990}). In
other words, the degree to which a worker was \emph{intrinsically
motivated} to do their work became a much more prominent focus for
interested researchers.

Drawing upon and contributing to the burgeoning psychology research on
motivation in the later half of the 1980s and early 1990s, a prominent
concept to emerge was \emph{psychological empowerment}, which was
theorized to be a multidimensional concept manifested as a set of four
distinguished `cognitions': \emph{meaning}, \emph{competence},
\emph{self-determination}, and \emph{impact}. The Psychological
Empowerment Scale was later developed, tested, and applied in multiple
studies inquiring into whether it served as a mediator for certain
workplace behaviors and turnover intention.

The concerted focus in the psychological empowerment literature on
individual differences between workers' feelings of empowerment,
however, isn't able to offer much in terms of how to improve
psychological empowerment among workers in general. Although some
recommendations are offered, it is difficult to recommend general
``one-size-fits-all'' solutions to increase psychological empowerment
for individual workers or an entire workforce. Seeing as how individual
people are quite different in terms of what motivates them, provides
meaning, makes them feel competent, and so on, the reviewed literature
is limited in that respect. As will be seen, for instance, one strategy
to improve empowerment is to change the way the workers interpret the
information gleaned from their workplace environment and outcomes
(\citeproc{ref-thomas1990}{Thomas and Velthouse 1990}).

The academic literature examining turnover intention provides a contrast
to the research focused on psychological empowerment. Research on
retention and turnover has tended to focus primarily on the workplace
context and organizational factors that influence a workers job
satisfaction and organizational commitment. Although empowerment is
often included as a factor within this line of research, the focus is
less concerned with individual level psychology and more concerned with
factors such as quality of pay, benefits, advancement opportunities and
the like. There are a few recommended intervention strategies to improve
retention, but pertain more toward shaping a friendly, inclusive, and
beneficial workplace environment.

This review summarizes a selection of the literature on the concept,
measurement, and examination of psychological empowerment as it
influences workplace behavior and worker attitudes. Also reviewed are
the interventions suggested which aim to reduce turnover, increase
empowerment, workplace satisfaction, and improve worker performance.

\section{Psychological Empowerment: theory and
measurement}\label{psychological-empowerment-theory-and-measurement}

\subsection{Powerlessness in the
workplace}\label{powerlessness-in-the-workplace}

Initially, the idea of \emph{powerlessness} in the workplace referred to
feeling a lack of autonomy and influence in the workplace
(\citeproc{ref-ashforth1989}{Ashforth 1989}). Autonomy was generally
understood as, ``The degree to which the job provides substantial
freedom, independence, and discretion to the employee in scheduling his
work and in determining the procedures to be used in carrying it out''
(\citeproc{ref-hackman1975}{Hackman and Oldham 1975}). Ashforth
(\citeproc{ref-ashforth1989}{1989}) argued that feelings of
powerlessness were incurred from unmet expectations and desires of
personal autonomy and involvement in the workplace. His argument was
that the greater the difference between expectations and perception of
reality in the workplace generates feelings of powerlessness. Such
powerlessness kicks off a contrived process of psychological adjustment
and coping. People first react to the difference between expectations
and reality, then develop increased feelings of inability to improve
their situation (i.e., helplessness) based on the realities of their
experience at work. Then, finally, come to feel detached and alienated
from the job and organization entirely.

\begin{quote}
``The findings, then, suggest that powerlessness and high generalized
expectations of control\footnote{Note that ``generalized expectations of
  control'' here simply refers to self-reported educational attainment
  on a 7-point scale.} give rise to unmet expectations/desires for
control; that the latter two constructs fuel reactance; that reactance
in turn, along with powerlessness, increases helplessness; and that
helplessness increases work alienation. The model suggests that the
experience of powerlessness can be quite traumatic, entailing
frustration, disruptive behaviors, feelings of helplessness, and loss of
job involvement and organizational identification. Conversely, it
suggests that fostering a sense of control can forestall frustration and
reactance, increase perceptions of organizational legitimacy and
self-efficacy, and stimulate involvement.''
(\citeproc{ref-ashforth1989}{Ashforth 1989, 234})
\end{quote}

\subsubsection{Recommended
Interventions}\label{recommended-interventions}

The recommended interventions of this early study were simple and
intuitive: make clear the actual link between the actions workers take
and the resulting outcomes they can expect to see in order to increase
perceptions of personal autonomy and control over workplace outcomes. To
this end, ``\ldots provide realistic job previews, specific and
challenging goals, prompt and balanced feedback, contingent
reinforcement, and appropriate role models\ldots{}'', and remove
obstacles which impede personal control
(\citeproc{ref-ashforth1989}{Ashforth 1989, 234}). In addition, it was
suggested that management relax or delegate control to those workers
whose jobs hold less actual influence or significance. By providing
opportunities for workers with ``low power jobs'' to participate in the
strategic and operational decisions usually confined to managerial
roles, workers may come to feel as though their input is not only valued
but influential upon the final decisions made\footnote{Note that this is
  the same reasoning explicated in procedural theories of representative
  democracy, though under different predicate assumptions
  (\citeproc{ref-dahl1961}{Dahl 1961}; \citeproc{ref-dahl2015}{Dahl and
  Shapiro 2015}).}. Although the reported findings of Ashforth
(\citeproc{ref-ashforth1989}{1989})'s study were narratively intuitive,
and the interventions sensible, it is unclear whether the theoretical
process of powerlessness results in powerlessness or begins from
it\footnote{From Ashforth's study, it's unclear exactly which variables
  are the dependent versus the independent variables, but it is clear
  that the primary variable of interest, powerlessness, is the
  independent variable of most importance to the author. Powerlessness,
  however, was measured by a combination of five different subscale
  measures derived from multiple sources, one of which included 32
  survey items. The measurement of powerlessness was assessed by a
  combination of 1) the 3-item Job Autonomy Scale from the Job
  Diagnostics Survey (\citeproc{ref-hackman1974}{Hackman and Oldham
  1974}, \citeproc{ref-hackman1975}{1975}), 2) the 4-item Job Authority
  scale from the revised Organizational Assessment Instrument, 3)
  self-reported job grade, 4) ratings of autonomy and participation
  (based on job title) by one job evaluation specialist, 5) a 32-item
  Powerlessness scale which consisted of a 15-item Participation
  subscale and a 17-item Autonomy subscale. It is not entirely clear
  where the 32-item powerlessness scale is adapted from.}.

\subsection{The Interpretive Model of
Empowerment}\label{the-interpretive-model-of-empowerment}

The difficulties and limitations of prior research prompted later
researchers to refine a more cohesive theoretical model of empowerment.
The publication by Thomas and Velthouse
(\citeproc{ref-thomas1990}{1990}) elaborates on a theoretical model of
empowerment based on review of relevant literature (up to the point of
publication in 1990)\footnote{Note that in the article cited, no test of
  the theoretical model is included, and no empirical results presented
  to validate the model. Instead, the authors rely on references to
  literature from psychology and clinical researchers. The abstract of
  the cited article mentions preliminary evidence that substantiates it,
  but the preliminary evidence mentioned in the abstract is confined to
  brief summaries of two different studies near the end of the article,
  where the evidence described suggests support for the authors
  interpretive model of empowerment. Thus, the review article argues for
  a particular non-traditional paradigm of management which fosters a
  worker's empowerment in theory.}. These authors conceptualized
empowerment as a multidimensional concept and operationalized
empowerment as \emph{intrinsic task motivation}, which refers to
positive experiences derived directly from doing some task. The
cognitive model developed involves identifying certain cognitions,
referred to as \emph{task assessments}, which manifests as a set of four
cognitions reflecting an individual's orientation to their task/work
role: \emph{meaning}, \emph{competence}, \emph{self-determination}, and
\emph{impact}. These task assessments are assumed to be the ``proximal
cause'' of intrinsic task motivation. That is to say, these four
``cognitions'' are considered latent constructs of their own, which in
turn manifest the second-order latent construct empowerment\footnote{The
  attempt to understand empowerment as ``intrinsic task motivation'' did
  not stick in the long run, but the multidimensional nature of the
  construct, as well as the four essential components, did.}.

An assumption of the conceptual definition of empowerment is that it is
not understood as an enduring personality trait generalizable across
situations. The set of four cognitions, i.e., task assessments, are said
to be shaped by the work environment. This fixes empowerment to the
specific workplace domain.

To put the theory in brief: People are empowered (i.e., energized) when
they are \emph{intrinsically motivated} to carry out some task, and they
feel \emph{intrinsically motivated} when they 1) judge the tasks as
sufficiently \emph{impactful} (i.e., the degree to which behavior is
seen as making a difference in terms of accomplishing the purpose of the
task), 2) judge themselves as \emph{competent} (i.e., self-efficacy), 3)
find the task \emph{meaningful}, and 4) feel as though their decision to
do the task is \emph{self-determined} (i.e., autonomous). They come to
judge specific tasks in this way by their particular kind of
\emph{interpretive style} and \emph{global beliefs} about work in
general as information is gleaned from event outcomes observed in the
workplace environment. That is to say, they interpret information from
the outcomes in whatever way is concordant to their interpretive style
and ``global assessments'' of tasks. When their interpretation of the
outcomes leads to higher task assessments (i.e., higher values of the
four cognitive variables), they then feel \emph{intrinsically} motivated
to carry out the task, resulting in behaviors in which they may
demonstrate \emph{flexibility}, \emph{initiative}, and
\emph{resilience}. Such behavioral outcomes are attributed purely to the
authors' theory on the effects of the four components of ``intrinsic
task motivation''.

Ultimately, this theory uses the word ``intrinsic'' when it is
unnecessary and not aligned with its meaning. The authors could remove
the word ``intrinsic'' entirely and it would make no difference; they
confuse what intrinsic motivation means and don't properly distinguish
it from extrinsic motivation. In so doing, this early work is undermined
by logical inconsistency. Supposing that a person's motivation to do
some task is indeed \emph{intrinsic}, then information on outcomes in
the environment should make little to no difference because the person
is motivated \emph{to do the task itself}. If the prospective outcome
drives the individual behavior, then the motivation is instrumental,
thus \emph{extrinsic}. There's a vast literature on the psychology of
motivation with multiple and competing perspectives, especially with
regard to notions of `intrinsic' v. `extrinsic' motives
(\citeproc{ref-deci2012}{Deci and Ryan 2012};
\citeproc{ref-ryan2000}{Ryan and Deci 2000},
\citeproc{ref-ryan2017}{2017}, \citeproc{ref-ryan2021}{2021}; for
contrast, see \citeproc{ref-kruglanski2017}{Kruglanski et al. 2017}; and
\citeproc{ref-kruglanski2022}{Kruglanski et al. 2022}).

Psychological theories vary considerably in explanations of behavior,
especially classics, whereas one theory maintains that all behaviors are
motivated by rewards (see \emph{operant theory},
\citeproc{ref-skinner1965}{Skinner 1965}), another asserts that all
behaviors are motivated by physiological drives, to which intrinsically
motivated behaviors were ones that provided satisfaction of
psychological needs (\citeproc{ref-hull1943}{Hull 1943}). More recent
explanations of behavior still vary considerably, although a relatively
recent Self-Determination Theory follows motivational theories to which
behavior seeks to satisfy psychological needs
(\citeproc{ref-ryan2000}{Ryan and Deci 2000}). In this work,
\emph{intrinsic motivation} refers to doing something because the
activity is inherently interesting or enjoyable itself; whereas doing
something because it leads to some separable outcome refers to
\emph{extrinsic motivation}. Although, the distinction isn't necessarily
as simple due to the fact that the classifying criteria for a behavior
to be motivated intrinsically (opposed to extrinsically) is quite high
and generally excludes many every-day behaviors. Ryan and Deci
(\citeproc{ref-ryan2000}{2000}), the progenitors of Self-Determination
Theory proposed that the need for competence underlies all motivation,
but full autonomy must be experienced for an activity to be truly
intrinsically motivated\footnote{Autonomy, in the framework of
  Self-Determination Theory, refers to the need to experience volition
  by controlling one's own behavior such that it integrates with one's
  sense of self, and the experience of autonomy depends on the degree to
  which a person's behavior is motivated by the level of external
  control versus the intrinsic motivation (for review see
  \citeproc{ref-higgins2021}{Higgins and Nakkawita 2021}).}.

\subsubsection{Recommended
Interventions}\label{recommended-interventions-1}

Two different kinds of interventions are thought to increase empowerment
(i.e., intrinsic task motivation): either adjust the working environment
outcomes (i.e., environmental changes) or re-frame how environment
outcomes are interpreted by individual workers (i.e., interpretive
intervention strategies). The \emph{environmental changes} intervention
strategy relates to the ideas of prior research which professes that
interventions modifying external variables such as organizational
structure, managerial strategies, and types of information all work as
interventions to produce increased feelings of empowerment. Thomas and
Velthouse (\citeproc{ref-thomas1990}{1990}) list a few environmental
variables which they theorized would shape the four cognitions/task
assessments: leadership, delegation, job design, and reward systems
(\citeproc{ref-thomas1990}{Thomas and Velthouse 1990}, see Table 1 on
page 676).

The latter ``\emph{interpretive intervention strategies}'' are aimed at
addressing/modifying how individual workers construe work outcomes and
consequences to increase worker empowerment. Only one ``interpretive
intervention strategy'' is mentioned as a feasible solution according to
the model: self-empowerment programs. Such programs are thought to help
individuals identify and practice styles theoretically conducive to
enhance tasks assessments crucial to increasing empowerment. The authors
say, ``The interpretive styles in the model are viewed as developed
\emph{habits} rather than innate abilities. These habits are relatively
isolated from day-to-day learning activities''
(\citeproc{ref-thomas1990}{Thomas and Velthouse 1990, 677}, emphasis
original). They go on to say,

\begin{quote}
``\ldots researchers/clinicians have demonstrated that such styles can
be changed by making the individual aware of assumptions that are
inherent in a style and by teaching individuals to consciously monitor
those ongoing interpretations and their consequences\ldots Thus,
according to the model,''self-empowerment'' programs are feasible
solutions to help individuals identify and practice styles of
attributing, evaluating, and envisioning, which would enhance their task
assessments.'' (\citeproc{ref-thomas1990}{Thomas and Velthouse 1990,
677})
\end{quote}

The authors don't detail what these self-empowerment programs look like,
but the goal is essentially to change the way the worker interprets the
things that happen at work in order to make such outcomes more
empowering. Offering support to the theoretical claims posited, the
authors cite to unpublished work\footnote{The evidence the authors cite
  to are two unpublished doctoral dissertations. The description of the
  two studies suggests that these dissertations were directly testing
  the cognitive model of empowerment they introduce in their review.
  However, these studies were unable to be located for direct citation.}
in which the four cognition task assessments of the empowerment model
mediated the relationship between positive feedback and self-reports of
intrinsic motivation. In other words, \emph{positive feedback}
supposedly caused increases in the four cognitive task assessments
(e.g., meaning, competence, self-determination, and impact) which, in
turn, increased self-reported ratings of intrinsic motivation
(\citeproc{ref-thomas1990}{Thomas and Velthouse 1990, 677}). However, it
is not clear how the four components of empowerment were measured.
Development of a survey measurement instrument didn't come about until
about five years after the model was theorized.

\subsection{Psychological Empowerment
Scale}\label{psychological-empowerment-scale}

The theoretical model existed for some time before an attempt was made
to develop a measure of psychological empowerment. Spreitzer
(\citeproc{ref-spreitzer1995}{1995}) developed and tested the validity
of a novel measure of psychological empowerment, following the cognitive
model of empowerment (i.e., ``intrinsic task motivation'') theorized by
Thomas and Velthouse (\citeproc{ref-thomas1990}{1990}) and based on
measurement scales adapted from previous research.

\begin{quote}
``In sum, psychological empowerment is defined as a motivational
construct manifested in four cognitions: meaning, competence,
self-determination, and impact. Together, these four cognitions reflect
an active, rather than a passive, orientation to a work role. By active
orientation, I mean an orientation in which an individual wishes and
feels able to shape his or her work role and context. The four
dimensions are argued to combine additively to create an overall
construct of psychological empowerment. In other words, the lack of any
single dimension will deflate, though not completely eliminate, the
overall degree of felt empowerment.''
(\citeproc{ref-spreitzer1995}{Spreitzer 1995, 1444})
\end{quote}

The scale includes 12-items with the component dimensions measured with
three items each. In this case, psychological empowerment is understood
as a common second-order latent factor that underlies the four latent
component factors (e.g., meaning, impact, competence,
self-determination)\footnote{See online appendix for additional notes on
  measurement for the scale.}.

\begin{longtable}[]{@{}
  >{\raggedright\arraybackslash}p{(\columnwidth - 2\tabcolsep) * \real{0.1944}}
  >{\raggedright\arraybackslash}p{(\columnwidth - 2\tabcolsep) * \real{0.8056}}@{}}
\caption{Psychological Empowerment 12-item Measurement scale
(\citeproc{ref-spreitzer1995}{Spreitzer
1995})}\label{tbl-empowerment}\tabularnewline
\toprule\noalign{}
\begin{minipage}[b]{\linewidth}\raggedright
Component
\end{minipage} & \begin{minipage}[b]{\linewidth}\raggedright
Item wording
\end{minipage} \\
\midrule\noalign{}
\endfirsthead
\toprule\noalign{}
\begin{minipage}[b]{\linewidth}\raggedright
Component
\end{minipage} & \begin{minipage}[b]{\linewidth}\raggedright
Item wording
\end{minipage} \\
\midrule\noalign{}
\endhead
\bottomrule\noalign{}
\endlastfoot
Meaning & \begin{minipage}[t]{\linewidth}\raggedright
\begin{itemize}
\tightlist
\item
  The work I do is very important to me.
\item
  My job activities are personally meaningful to me.
\item
  The work I do is meaningful to me.
\end{itemize}
\end{minipage} \\
Competence & \begin{minipage}[t]{\linewidth}\raggedright
\begin{itemize}
\tightlist
\item
  I am confident about my ability to do my job.
\item
  I am self-assured about my capabilities to perform my work activities.
\item
  I have mastered the skills necessary for my job.
\end{itemize}
\end{minipage} \\
Self-Determination & \begin{minipage}[t]{\linewidth}\raggedright
\begin{itemize}
\tightlist
\item
  I have significant autonomy in determining how I do my job.
\item
  I can decide on my own how to go about doing my work.
\item
  I have considerable opportunity for independence and freedom in how I
  do my job.
\end{itemize}
\end{minipage} \\
Impact & \begin{minipage}[t]{\linewidth}\raggedright
\begin{itemize}
\tightlist
\item
  My impact on what happens in my department is large.
\item
  I have a great deal of control over what hap pens in my department.
\item
  I have significant influence over what happens in my department.
\end{itemize}
\end{minipage} \\
\end{longtable}

Overall, the development and test for the original measurement scale for
psychological empowerment can be viewed as a first step towards
measuring empowerment in a work related context. However, no real useful
measurement instrument for psychological empowerment is made readily
available from this study. The methods used to test validity and
reliability of the four measurement scales for empowerment is
interesting but complex, and doesn't appear to result in an actual
estimate of the latent construct for empowerment itself. Instead,
convergent validity is obtained but not discriminant validity for the
four components of empowerment. Interested researchers might be better
off developing and validating a single measurement scale for empowerment
itself rather than rely on measuring the presence of four separate
latent component constructs with 12 requisite survey items.

\section{Psychological Empowerment as a
Mediator}\label{psychological-empowerment-as-a-mediator}

Multiple subsequent studies inquired into whether psychological
empowerment functioned as a mediator for workplace performance, or
incorporated the concept or measurement thereof into studies
investigating the antecedents to turnover intentions. This line of
research generally finds that psychological empowerment plays a positive
mediating role in helping to increase Organizational Citizenship
Behavior (OCB)\footnote{Organizational Citizenship Behavior is defined
  differently by different scholars, but is best understood as voluntary
  work behavior done within an organization that is not required as part
  of the worker's duties of employment. It is a blanket term that refers
  to worker behavior that contributes benefit to the organization but is
  beyond their formal obligations.} and levels of performance
(\citeproc{ref-taylor2013}{Taylor 2013}), workplace satisfaction and
organizational commitment (\citeproc{ref-liden2000}{Liden, Wayne, and
Sparrowe 2000}), and mediates the effect of perception of authentic
leadership from one's supervisor
(\citeproc{ref-shapira-lishchinsky2014}{Shapira-Lishchinsky and Tsemach
2014}). However, the latter two of the cited studies
(\citeproc{ref-liden2000}{Liden, Wayne, and Sparrowe 2000};
\citeproc{ref-shapira-lishchinsky2014}{Shapira-Lishchinsky and Tsemach
2014}), which explore psychological empowerment as mediator for other
outcomes, are questionable and should be reviewed with increased
scrutiny before the reported findings are readily accepted\footnote{See
  online appendix for expanded scrutiny of the latter two articles
  cited.}.

Research by Taylor (\citeproc{ref-taylor2013}{2013}) is the least
problematic of these studies and is easily summarized. Goal setting
plans by organizations positively influenced psychological empowerment
and helpful and cooperative behavior toward other employees in the
organization (i.e., `organizational citizenship behavior' toward
individuals). Goal setting increased self-reports of individual-oriented
organizational performance behavior (OCBI)\footnote{OCBI : ``I help
  others who have heavy workloads'', ``I take time to listen to
  co-workers' problems and worries''. ``I go out of my way to help new
  employees.'' ``I assist my supervisor with his/her work even when I am
  not asked.''} more so than self-reported organization-oriented
performance behaviors (OCBO)\footnote{OCBO: ``My attendance at work is
  above the norm'', ``I obey informal rules that maintain order'', ``I
  give advance notice when I am unable to work''}. Perceived difficulty
of goals\footnote{Goal-difficulty: ``A high degree of skill and know-how
  is necessary to do my job well.'', ``Jobs like mine are quite
  demanding day after day.''} positively influenced both OCBI and OCBO.
Perception that job goals are clear, specific, and fully understood
(goal-specificity)\footnote{Goal-specificity: ``My responsibilities at
  work are very clear and specific.'', ``I understand fully which of my
  job duties are more important than others.''} did not directly
influence either OCBI or OCBO. However, goal-specificity indirectly
influenced OCBI through the indirect effect of psychological
empowerment. In other words, perception that job goals are clear,
specific, and fully understood positively influences psychological
empowerment which in turn indirectly influences individual-oriented
performance (OCBI). Goal-specificity and goal-difficulty both positively
influenced psychological empowerment, but the effect of goal-specificity
was larger in comparison.

\subsubsection{Recommended Intervention
Strategies}\label{recommended-intervention-strategies}

Review of the aforementioned studies conveys the impression that
anything which may increase any of the four components of psychological
empowerment should result in positive improvement in worker performance,
enthusiasm on the job, and helpful workplace behaviors. This, however,
tells us little more than what may be obvious without providing specific
interventions on how to go about doing so. That is, the results of
Taylor (\citeproc{ref-taylor2013}{2013}) are derived from comparison of
attitudinal surveys but stops short of testing specific intervention
strategies. Moreover, the sample was comprised of Australian public
service workers, which may limit external validity--not to suggest
anything special about Australia but rather the potential differences
between public service workers and others generally. Further research
may help to confirm whether interventions providing specific and
challenging goals promotes improved worker performance, organizational
commitment, or other such Organizational Citizenship Behaviors.

Perceiving the boss to be authentic, or otherwise nice, good, and
honest, increases the degree to which a worker feels like they have an
impact at work, but the research showing this also shows that perceiving
the boss to be authentic also increases one's willingness to find work
elsewhere (\citeproc{ref-shapira-lishchinsky2014}{Shapira-Lishchinsky
and Tsemach 2014}). Said another way, perceiving the boss to be an
authentic leader \footnote{``Authentic leadership'' was measured using
  four survey items: ``My principal says exactly what he or she means'',
  ``My principal admits to mistakes when they are made'', ``My principal
  tells me the hard truth'', and ``My principal outwardly displays how
  she/he feels''. (See Appendix in
  \citeproc{ref-shapira-lishchinsky2014}{Shapira-Lishchinsky and Tsemach
  2014})} \emph{increases} a worker's intent to leave their current job.
The recommendation when it comes to leadership is to ``\ldots promote
high standards of authentic leadership''
(\citeproc{ref-shapira-lishchinsky2014}{Shapira-Lishchinsky and Tsemach
2014, 701}), which is not a recommendation but has more to do with a
worker's perception of the boss's personal characteristics. Deriving
meaning and feeling competent in one's work significantly influences
work satisfaction, organizational commitment, and better ratings of job
performance from one's supervisor (\citeproc{ref-liden2000}{Liden,
Wayne, and Sparrowe 2000}), but this research is highly dubious since it
shows findings that are contradictory: deriving meaning from one's work
increases work satisfaction but feeling competent in one's work also
\emph{decreases} work satisfaction. These perplexing findings are likely
the result of limitations or flaws of the research. Attempts to
replicate findings in the future may help our understanding in this
area.

The best intervention suggestion comes from Taylor
(\citeproc{ref-taylor2013}{2013}),

\begin{quote}
``The findings suggest that public managers can promote OCB from
goal-setting plans in two ways. One is to ensure that the goal-setting
plans are well designed and implemented. In particular, challenging
goals that are accepted by employees can motivate higher levels of
performance. Another is to improve employees' views of their work
environment and their individual capacity so that they feel
psychologically empowered at work. To encourage employees to see their
job as more meaningful and to satisfy their self-efficacy and
self-determination needs, managers should help employees understand the
scope of their job and responsibilities, articulate the overlap between
the agency's goals and the employees' personal goals, and demonstrate
support for employees in their pursuit of the goals (e.g., provide
access to appropriate professional development programs)''
(\citeproc{ref-taylor2013}{2013, 461}).
\end{quote}

Thus, organizations that set specific and difficult job goals may
increase workers psychological empowerment and helpful performance
behavior towards individuals in the organization. In particular, setting
job goals may positively influence psychological empowerment which in
turn is expected to improve levels of performance directed to help
co-workers and benefit the broader organization. Psychological
empowerment increases for employees who feel their job goals require a
high degree of competence and are demanding. The same goes for employees
who feel that their work goals are clear, specific, and fully
understood.

\section{Turnover Intention}\label{turnover-intention}

A related line of research in management and organizational studies
inquired into factors related to voluntary turnover, which include
empowerment among others such as job satisfaction and organizational
commitment. However, these studies do not maintain a strict focus on the
presence and effect of psychological empowerment, and the 12-item scale
is not used to measure empowerment in the following studies. At least
one study measures empowerment by use of a single item measure on a
5-point agree/disagree Likert response option format, ``Employees have a
feeling of personal empowerment with respect to work processes.''
(\citeproc{ref-pitts2011}{Pitts, Marvel, and Fernandez 2011}). Other
research includes measurement of constructs which tangentially relate to
the conceptual components of psychological empowerment without
referencing them as such. For instance, as part of the measure for
``functional preferences'' Bertelli (\citeproc{ref-bertelli2007}{2007})
measures ``intrinsic motivation'' with three survey items drawn from the
Federal Human Capital Survey (FHCS): ``Employees have a feeling of
empowerment and ownership of work processes.'', ``My work gives me a
feeling of personal accomplishment.'', ``I like the kind of work I do.''
The other half of the comprised ``functional preferences'' measure
includes items intent on measuring \emph{job involvement}: ``My talents
are used well in the workplace.'', ``I recommend my organization as a
good place to work.'', ``The work I do is important.''

The variable use of different measurement items speaks to the
discrepancies noted in results across multiple studies by Tett and Meyer
(\citeproc{ref-tett1993}{1993}). In their meta-analysis, they attempt to
untangle the variable findings from the disparate use of multiple
measures of similar or identical constructs. Their aim was to compare
job satisfaction and organizational commitment ``as unique precursors of
employee withdrawal\ldots{}''

\begin{quote}
``Overall, results support the view that satisfaction and commitment
each contribute uniquely to the turnover process. Those contributions,
however, are largely limited to intention/cognitions, and depend on the
choice of measures.'' (\citeproc{ref-tett1993}{Tett and Meyer 1993,
279}).
\end{quote}

The authors say that job satisfaction and {[}self-reported{]} turnover
intention is sometimes included into an index of ``withdrawal
cognitions'', which refers simply to thoughts of quitting and intention
to search for alternative employment. However, turnover intention is
conceived as a conscious and deliberate willfulness to leave the
organization. That is, turnover intention is conceptually different from
a set of ``withdrawal cognitions''. In other words, definitive intention
to leave the job is different in that it is serious and intentional
rather than merely passive thoughts about quitting.

The first set of hypotheses (Hypothesis 1a and 1b) expect that
organizational commitment should show stronger negative correlation than
job satisfaction with turnover intention (H\textsubscript{1a}) and
actual turnover (H\textsubscript{1b}). The authors found that job
satisfaction showed a stronger negative relationship with turnover
intention than when compared with organizational commitment. This would
imply, simply, that higher job satisfaction is negatively correlated
with turnover intention. However, once the more deliberate and serious
turnover intention is detached from the aforementioned set of passive
``withdrawal cognitions'', the strength of the negative relationship is
lessened substantially. Ultimately, the findings suggest that
organizational commitment and job satisfaction are relatively equal in
importance for predicting either serious turnover intention or
withdrawal cognitions (\citeproc{ref-tett1993}{Tett and Meyer 1993,
280}).

What is important to note, however, is that organizational commitment
showed a much stronger negative relationship with actual turnover
(contra. Passive thoughts or deliberate intention). The point is not
only that organization commitment is likely more important than job
satisfaction at predicting actual turnover, but crucially that an
``intention to leave'' measure is not viable as a proxy for actual
turnover despite its predictive power (\citeproc{ref-tett1993}{Tett and
Meyer 1993}). Perhaps just as important, the authors also note concerns
with available measures for job satisfaction. They question the content
of job satisfaction measures used in prior studies owing to
``\ldots contamination by items pertaining to withdrawal cognitions''
and present a list of studies where such contamination was evident
(\citeproc{ref-tett1993}{Tett and Meyer 1993, 283}). In other words,
items measuring job satisfaction commonly include items that also
measured thoughts about quitting.

The ultimate takeaway appreciable here is that researchers interested in
ascertaining whether job satisfaction or organizational commitment are
related to turnover are probably better off constructing and validating
measurement scales \emph{ex novo} (i.e., from the beginning).

Beyond the accounting of the study by Tett and Meyer
(\citeproc{ref-tett1993}{1993}), this review finds that scholars from
different fields approach questions related to empowerment in the
workplace, job satisfaction, turnover and intention, and organizational
commitment in distinct ways that place stronger emphasis on some factors
over others. Of course, this is obvious--a trained economist is apt to
approach questions about turnover intent by way of reviewing and
comparing an organization's incentive structure. This seems to be the
case for the study conducted by Bertelli
(\citeproc{ref-bertelli2007}{2007}), who compared the effect of the
aforementioned ``functional preferences'' as a determinant for turnover
intention between two distinct samples--one where performance-based pay
reforms were put in-place for supervisors but not for non-supervisors,
compared to another where no such merit-based pay reforms were enacted.

To summarize, Bertelli (\citeproc{ref-bertelli2007}{2007}) finds that it
is easier to retain supervisors in federal agencies (the Internal
Revenue Service, IRS, and Office of the Comptroller of the Currency,
OCC) who feel like their work is important, who like their work and gain
feelings of accomplishment from it (i.e., ``functional preferences'').
This applies to employees who are Non-supervisors too, but
non-supervisors are also more likely to stay when they view their
workplace as friendly. A friendly work environment is not related to
turnover intention for supervisors. Perceptions of inadequate
pay\footnote{Q58 from the Federal Human Capital Survey (FHCS): ``How do
  you rate the amount of pay you get on your job?'' (1=very
  poor\ldots5=very good).} impacts the likelihood of turnover intention
for non-supervisors. To put simply, federal agency jobs like the IRS and
OCC for people in non-supervisory positions that are not perceived as
financially and psychologically rewarding are more likely to see
increased turnover from non-supervisors. Having friendly coworkers
helps, as does feeling like the work is important and feelings of
accomplishment, but perceptions of inadequate pay, micromanaging boss,
and merit-based performance incentives hanging over head all work to
increase likelihood for turnover among non-supervisors. Those in
supervisor positions are best retained when they find their work
important and personally rewarding--perceptions of inadequate pay
doesn't drive turnover intention for supervisors.

A brief note on the method. Bertelli writes
(\citeproc{ref-bertelli2007}{2007}),

\begin{quote}
``The OCC is included in this study as a''control'' group for examining
the effect of perceptions of the management incentives unique to the
IRS. OCC, located with the IRS in the U.S. Department of the Treasury,
is the regulator of U.S. banks, and its employees engage in similar
accounting oversight (i.e., auditing) as that performed by the IRS at
the bank rather than at the individual and corporate levels. Both
subagencies have a nationwide staff, and their employees possess
generally similar backgrounds in accounting and financial matters. If
there are unobserved attributes correlated with the responses to
questions regarding performance incentives specific to the IRS, they can
reasonably be addressed by the comparison between results for IRS and
OCC supervisors. Among those attributes is the pay-banding system. In
the absence of questions related specifically to pay-banding to
facilitate a direct test, comparison of these groups is suggestive
evidence of the effect of pay-banding on bureaucrats' intent to leave
their jobs. When combined with the information included in survey
responses, this comparison provides significant leverage in
understanding the effects of performance incentives on turnover
intention.'' (\citeproc{ref-bertelli2007}{2007, 239--40})
\end{quote}

The results of a Probit regression analysis for determinants of turnover
intention for supervisors across the two samples shows that only
``functional preferences'' and being employed by the IRS negatively
predict turnover intention. Bertelli's study takes the difference of
agency employment as a proxy for the effect of the pay-banding reform
scheme--that is, ``In the absence of questions related specifically to
pay-banding to facilitate a direct test, comparison of these groups is
suggestive evidence of the effect of pay-banding on bureaucrats' intent
to leave their jobs.'' (\citeproc{ref-bertelli2007}{Bertelli 2007,
240}).

An alternative interpretation may omit any acknowledgement of the
pay-banding scheme enacted for supervisors in the IRS, stating simply be
that working for the OCC as a supervisor positively predicts turnover
intention among supervisors. The author relies on the strong assumption
that the two samples match closely enough that differences attributed to
agency employment serve to indicate the effect of the performance-based
pay reforms. However, by this logic, any differences between the two
agencies in organizational or employee policies pertaining to
supervisors can account for the differential effect. It isn't simply
that the model doesn't account for every possible difference between the
two agencies, but rather that the dummy effect of being employed as a
supervisor in one agency rather another is posited to indicate the
effect of this specific pay-banding scheme enacted within the IRS.

Additionally, Bertelli (\citeproc{ref-bertelli2007}{2007}) notes the
striking insignificance of negative perceptions of pay quality for
supervisors in both the IRS and OCC samples. A single-item question
measures this attitude about quality of pay, ``How do you rate the
amount of pay you get on your job?'' where response options range across
a 5-point scale from very poor to very good. From what can be gathered,
it appears that higher ratings about quality of pay were not significant
predictors of turnover intention. The mean response for quality of pay
across the whole sample (\(n = 1842\)) was \(3.635\) with a standard
deviation of \(1.003\). The mean response for supervisors only was not
displayed in the study.

As an aside, Bertelli measures quality of pay with a single-item, but
does not take into consideration any other forms of compensation,
especially those unique benefits related to federal employment compared
to similar jobs in the private sector. Anecdotally, private contractors
recruit large amount of veteran service members for their specific
service-related skill sets by offering appealing payment incentives; but
working as a private contractor does not always come with the benefits
or other compensatory incentives afforded by work as an employee, nor as
an employee of the federal government. The decision between working in
the public or private sector for the same kind of work may come down to
such a comparative assessment of the financial security offered by
federal government employment over self-employed contractor work.
Although the compensation across the two agencies is likely to be
similar, if not identical, the notion that quality of pay or
compensation/benefits is insignificant to turnover intention, or actual
turnover for that matter, is questionable. Especially considering that
pay quality across both agencies was a significant predictor of turnover
intention for non-supervisors. Bertelli doesn't include descriptions of
the average reported income for the studied samples, nor a reference to
the relative pay scales for government employees that would apply to
either agency, nor is any consideration made with regard to the
overarching compensation/benefits structure. The quality of pay is
demonstrably insignificant for turnover intention, but this should not
be striking nor surprising for those who have reached supervisor status
while employed with the federal government; they are paid quite well
overall.

In contrast, Pitts, Marvel, and Fernandez
(\citeproc{ref-pitts2011}{2011}) does take \emph{benefits satisfaction}
into account when inquiring why federal employees choose to leave either
their respective agencies or government work outright\footnote{The
  dependent variables--turnout intention to leave agency and turnout
  intention to leave government--are derived from the same survey item
  question, ``Are you considering leaving your organization within the
  next year, and if so, why?'' How the authors are able to create two
  variables from one question with a dichotomous response of yes or no
  is beyond me. They don't state that the Federal Human Capital Survey
  in 2006 asks the question twice for both agency and government
  turnover intentions, so the reader is left to presume such is the
  case.}. However, these authors also minimize the impact of pay on
turnover intention but rely on self-reports of pay
satisfaction\footnote{Pay satisfaction was measured with a single-item
  on a 5-point Likert scale: ``Considering everything, how satisfied are
  you with your pay?''} while also asking confounding satisfaction
questions related to benefits (e.g., health insurance, retirement
benefits, life insurance)\footnote{Benefits satisfaction was measured
  with three items on a 5-point Likert scale: ``How satisfied are you
  with retirement benefits?'', ``How satisfied are you with health
  insurance benefits?'', ``How satisfied are you with life insurance
  benefits?''}, opportunities for advancement\footnote{Advancement
  Opportunity Satisfaction was measured with a single-item on a 5-point
  Likert scale: ``How satisfied are you with your opportunity to get a
  better job in your organization?''}, and even job
satisfaction\footnote{Job satisfaction was measured with a single-item
  on a 5-point Likert scale: ``Considering everything, how satisfied are
  you with your job?''}. All these are arguably related to the same
underlying factor of compensation someone takes into consideration when
searching for a job. And this notional latent factor is likely crucial
to understand the reasons for both recruitment \emph{and} retention. The
reasons why a person intends to leave may not be related to their
financial situation if said situation is far more comfortable in
government work than elsewhere. Yet the reason for choosing to get into
government work in the first place is likely strongly related to the
perceived prospective financial security of working in the federal
government, and might be a foundational reason for the prevalence of
inter-agency turnover compared to government turnover in general. About
\(16.1\%\) of federal employees in the sample stated their intent to
leave their agency for another agency within the federal government. In
comparison, about \(3.6\%\) of federal employees plan to leave
government work all together\footnote{Pitts, Marvel, and Fernandez
  (\citeproc{ref-pitts2011}{2011}) don't specify the number of
  observations explicitly. Rather, the number of observations in the
  survey is found in Table 2 on page 755 (\(n = 217,504\)). From the
  cumulative percentages of those indicating turnover, whether leaving
  agency or government, that amounts to about \(42,848\) employees from
  the sample who reported an intention to leave within the next year
  (assuming my math is correct). The count of those who reported
  intention to leave an agency but stay working within federal
  government is approximately \(35,018\). Whereas the approximate count
  of those who reported intention to leave federal government work
  overall is \(7,830\). Seeing as how leaving government is far less
  substantial a figure than leaving an agency, I see the reason for the
  distinction between the two but it would have been better to include
  an actual turnover rate of some sort for the past year at least. A
  statistic from the Bureau of Labor Statistics, or some other
  concurrent figure that demonstrates just how significant turnover
  actually is for government employees. The authors mention some
  consistent 20 year decline, but don't present the evidence of that.}.

\subsubsection{Recommended Interventions to Reduce Turnover and Turnover
Intention}\label{recommended-interventions-to-reduce-turnover-and-turnover-intention}

The literature reviewed here is mostly concerned with turnover intention
for government employees. To summarize the findings and recommended
interventions of Pitts, Marvel, and Fernandez
(\citeproc{ref-pitts2011}{2011}), workplace satisfaction plays largest
role in predicting intention to leave an agency among government
employees. Advancement opportunities and age also plays a significant
role for agency turnover intention. The most important factors
predicting whether an employee intends to leave federal government are
age, whether an employee is non-White, and job satisfaction.

These authors suggested a couple of interventions to decrease turnover
for government agencies. First, offer advancement opportunities along
with professional development programs in hopes that it will motivate
new and younger employees to say at their agency instead of move between
government agencies. Second, relax performance-based workplace culture
(merit-based pay/advancement, incentives) for new and younger employees
as such performance-based cultural work environment predicts
intra-agency turnover intention among this group
(\citeproc{ref-pitts2011}{Pitts, Marvel, and Fernandez 2011}).

As stated previously, Bertelli (\citeproc{ref-bertelli2007}{2007}) finds
that ``functional preferences'' (i.e., a latent factor indicated by two
component measurements of job involvement and intrinsic motivation) and
solidarity preferences (e.g., ``This is a friendly place to work.'')
decrease the likelihood that agency bureaucrats will leave their
government jobs (\citeproc{ref-bertelli2007}{2007, 253}). Bertelli does
not recommend interventions outright, there may be distinct strategies
implied depending on whether the target workers in question are
supervisors or non-supervisors. However, since the perception of a
friendly work environment reduces the likelihood of turnover intention
among non-supervisor employees, it is unclear how or what interventions
could be enacted to make the workplace appear friendly, generally
speaking. Although performance-based incentives are weakly significant
in predicting whether supervisors retain their jobs, such incentives
imposed upon non-supervisor employees may increase turnover intention.
However, bear in mind that turnover intention may not reflect actual
turnover rates (\citeproc{ref-tett1993}{Tett and Meyer 1993}).

\section{Discussion and Conclusion}\label{discussion-and-conclusion}

Based on the review of the cited literature related to turnover
intention for government employees, the findings suggest that government
employees are generally more likely to stay for obvious reasons such as
financial security and retirement potential. It seems more apparent that
simply getting a federal government job matters more than which agency
one settles with, and the longer one stays working in government then
the less likely they are to leave. For instance, one finding suggests
that a workplace environment with a ``performance culture'' (i.e., a
work environment that incorporates merit-based pay and advancement
incentives) \emph{increases} turnover intention but only inter-agency
turnover, i.e., those who plan to leave their agency to take another job
within the federal government. Absent here is consideration as to
whether government work in general is perceived to be more attractive
compared to work outside of the federal government for most current
government employees. If it is, then the reasons for wanting to leave
are unlikely to be entirely related to income and benefits. That is to
say, there is likely a substantive difference between those who work for
federal government agencies and those who don't.

This references back to the earlier literature on psychological
empowerment, which emphasizes a focus on the individual differences
between workers in terms of the meaning, impact, competence, and
autonomy they derive from the work they do. This literature draws from
organizational psychologically and recommends interventions related to
re-framing the workers interpretation of workplace outcomes and proposes
the increased use of self-empowerment programs
(\citeproc{ref-thomas1990}{Thomas and Velthouse 1990}), and generally
recommends setting goals that are both challenging and well understood
(i.e., specific) (\citeproc{ref-taylor2013}{Taylor 2013}). Although the
empowerment literature makes it clear that work satisfaction and
positive perceptions of workplace leadership either increase or are
increased by psychological empowerment, there is no clear way to
intervene and adjust the perceptions of the individual worker. In that
respect, it appears that the early recommendations by Ashforth
(\citeproc{ref-ashforth1989}{1989}) may be the most practical and
sensible overall, ``\ldots provide realistic job previews, specific and
challenging goals, prompt and balanced feedback, contingent
reinforcement, and appropriate role models\ldots{}'', and remove
obstacles which impede personal control
(\citeproc{ref-ashforth1989}{1989, 234}). In addition, delegating more
influence to lower level employees in ways that make their input
significant to workplace strategies and decision-making may also
increase some demonstrably important component factors of empowerment
such as competence, impact, and self-determination (i.e., autonomy).

In contrast, the turnover literature reviewed here focuses primarily on
the organizational structure and workplace environment intervention.
Although workers who feel like their work is important and gain feelings
of accomplishment from it are far easier to retain, these attitudes
become more significant to those who are in supervisor roles--i.e.,
those who are arguably already well settled into a job or career. The
workplace environment and its overarching structural arrangement is most
important to reduce turnover among the young and the restless
(\citeproc{ref-bertelli2007}{Bertelli 2007}). Quality of pay, benefits,
opportunities for advancement, and overall financial security are
organizational factors that are significant to consider if the goal is
to reduce turnover intention (\citeproc{ref-pitts2011}{Pitts, Marvel,
and Fernandez 2011}), and possibly actual turnover rates as well.
Although recruiters and hiring managers might love to hear a potential
hire genuinely respond to the question ``Why do you want to work here?''
with ``To increase my sense of psychological empowerment'', the real
motivation for many potential and current employees may still be those
basic instrumental factors.

A major aspect underlying the literature reviewed here is organizational
contextual factors versus individual differences. Thomas and Velthouse
(\citeproc{ref-thomas1990}{1990}) contended that empowerment refers to
the motivational content of a nontraditional form of management (a new
paradigm of motivation, as they say) in which management establishes
controls to foster an emphasis or internalized commitment to the task
\emph{itself}. Workers are ``pulled'' toward doing the tasks rather than
``pushed'' towards it because the tasks are made to be meaningful;
workers identify with the task; and there is expressive value in doing
the task. To put another way, figuring out ways to have workers want to
do the tasks rather than simply making them do the tasks. Authors of the
empowerment literature say that prior research professes the idea that
external variables such as organizational structure, managerial
strategies, and types of information all work as interventions to
produce increased feelings of empowerment (i.e., self-efficacy) in work.
In other words, the structural context of the organized work environment
induces greater empowerment within the worker, and thus enhances their
motivation to initiate, persist, and better perform. The empowerment
authors view this as lacking as it doesn't take into account how the
perspective of the subject (i.e., subjective) influences the cognitive
variables thought to determine feelings of empowerment. Stated
otherwise, prior research doesn't take into account the influence of
individual differences (\citeproc{ref-thomas1990}{Thomas and Velthouse
1990}).

Yet it appears as though that when organizational factors are relatively
stable and insignificant--mainly for those workers who are in
comfortable and secure supervisory positions--then turnover intent and
actual turnover rates among them are likely to be more strongly
associated with individual differences such as the components of
psychological empowerment. Those who are relatively more secure and in
roles in which they not only feel, but actually are in fact, impactful
and competent, are seemingly more likely to appraise their work
conditions in terms of how satisfying it is to them, how meaningful it
is, and how intrinsically motivated they are to continue doing such
work. In essence, when certain needs are sufficiently met, then other
factors take on more significance. This may explain why federal
employees are more inclined to change agencies but do not quit
government work outright (\citeproc{ref-bertelli2007}{Bertelli 2007}).
In contrast, psychological empowerment and other such factors are muted
\emph{unless} the individual worker is in fact legitimately
\emph{intrinsically motivated} to do that work for the sake of doing the
work \emph{itself}. Otherwise, it should be presumed that most workers
retain jobs motivated by external factors related to basic instrumental
incentives. It may be the case that workers seek out particular jobs or
types of work based on more \emph{internalized and integrated} extrinsic
motivations, which are the closest approximations to intrinsic
motivation along the extrinsic to intrinsic motivational spectrum
(\citeproc{ref-ryan2000}{Ryan and Deci 2000},
\citeproc{ref-ryan2017}{2017}). That is to say, the initial choice of
what work one wishes to get into is likely more powerful than palliative
interventions such as self-empowerment programs.

\newpage

\section{Appendix A}\label{appendix-a}

\subsection{The Public Service Motivation Construct and
Scale}\label{the-public-service-motivation-construct-and-scale}

The following is a brief overview of the Public Service Motivation scale
and construct as developed by Perry (\citeproc{ref-perry1996}{1996}).
Perry defines Public Service Motivation (PSM), ``as an individual's
predisposition to respond to motives grounded primarily or uniquely in
public institutions'' (\citeproc{ref-perry1996}{1996, 6}). An
interesting aspect of the PSM construct and scale is that it is
predicated on a particular understanding of motivation rooted in
evolutionary psychology, and is also found elsewhere in social
psychology literature (For review, see
\citeproc{ref-higgins2021}{Higgins and Nakkawita 2021, 88}). Here and in
other related work, motives are driven by needs. In particular, certain
``drive theories'' of motivation suggest that humans recognize something
as a `need' when its deprivation threatens our survival, our growth, or
our ultimate life satisfaction and well-being. The motivation to fulfill
unmet `needs' may guide our behavior, impact on our thoughts, explain
our selective attention and memory (\citeproc{ref-gardner2000}{Gardner,
Pickett, and Brewer 2000}; \citeproc{ref-hull1943}{Hull 1943};
\citeproc{ref-mcleod2022}{Mcleod 2022}). As it relates to PSM, Perry
(\citeproc{ref-perry1996}{1996}) recognized three analytically distinct
motivational categories: rational, norm-based, and affective.

\begin{quote}
``Rational motives involve actions grounded in individual utility
maximization. Norm-based motives refer to actions generated by efforts
to conform to norms. Affective motives refer to those triggers of
behavior that are grounded in emotional responses to various social
contexts.'' (\citeproc{ref-perry1996}{Perry 1996, 6})
\end{quote}

The PSM scale development was based on conceptual dimensions which were
categorized into the motivational categories. The final PSM scale was
confirmed as having four component dimensions: attraction to public
policy making, commitment to the public interest, compassion, and
self-sacrifice. However, the initial analysis of the finalized PSM scale
demonstrated no distinct difference in model fit between a three and
four dimensional PSM model. The public interest and self-sacrifice
component dimensions were highly correlated, and so a comparison between
the three and four-dimensional model was done. The four dimensional
model showed a slighter higher \(𝛘^2\) statistic, thus Perry
(\citeproc{ref-perry1996}{1996}) opted to retain the self-sacrifice
dimension to maintain the four dimensional model of the PSM scale
because, ``it has had a historical connection to how we think about
public service that is explicitly preserved by retaining the
dimension.'' (\citeproc{ref-perry1996}{Perry 1996, 20}). However, it is
recommended by Perry that further refinements to the PSM construct and
scale be undertaken by future research.

The four component dimensions of the Public Service Motivation scale are
described below;

\begin{description}
\item[Attraction to public policy making (rational motive)]
Refers to a rational motive in which the opportunity to participation in
the formation of public policy draws individuals to public service
\item[Commitment to the public interest (norm-based)]
An altruistic desire to serve the public interest
\item[Compassion (aka ``patriotism of benevolence'', affective)]
An affective state that combines the love of regime values and the love
of others which drives people to exemplify ``a patriotism of
benevolence''{[}\^{}22{]}
\item[Self-sacrifice (unclear motive)]
Refers to the willingness to substitute service to others for tangible
personal rewards. In other words, the willingness of public servants to
forego financial rewards for the intangible rewards received from
serving the public.
\end{description}

After refinements were made to the initial scale, the resulting Public
Service Motivation Scale was finalized to include 24-items measuring
four component dimensions. The PSM measurement items are measured with
Likert type response options. The items are listed in the table below

\begin{longtable}[]{@{}
  >{\raggedright\arraybackslash}p{(\columnwidth - 4\tabcolsep) * \real{0.1015}}
  >{\raggedright\arraybackslash}p{(\columnwidth - 4\tabcolsep) * \real{0.2386}}
  >{\raggedright\arraybackslash}p{(\columnwidth - 4\tabcolsep) * \real{0.6548}}@{}}
\caption{Public Service Motivation Scale Item Wording
(\citeproc{ref-perry1996}{Perry 1996})}\tabularnewline
\toprule\noalign{}
\begin{minipage}[b]{\linewidth}\raggedright
Motive type
\end{minipage} & \begin{minipage}[b]{\linewidth}\raggedright
Component Dimension
\end{minipage} & \begin{minipage}[b]{\linewidth}\raggedright
Item Wording
\end{minipage} \\
\midrule\noalign{}
\endfirsthead
\toprule\noalign{}
\begin{minipage}[b]{\linewidth}\raggedright
Motive type
\end{minipage} & \begin{minipage}[b]{\linewidth}\raggedright
Component Dimension
\end{minipage} & \begin{minipage}[b]{\linewidth}\raggedright
Item Wording
\end{minipage} \\
\midrule\noalign{}
\endhead
\bottomrule\noalign{}
\endlastfoot
Rational Motive & Public Policy Making &
\begin{minipage}[t]{\linewidth}\raggedright
\begin{itemize}
\tightlist
\item
  PSM 11 Politics is a dirty word. (Reversed)
\item
  PSM 27 The give and take of public policy making doesn't appeal to me.
  (Reversed)
\item
  PSM 31 I don't care much for politicians. (Reversed)
\end{itemize}
\end{minipage} \\
Norm-based motive & Public Interest &
\begin{minipage}[t]{\linewidth}\raggedright
\begin{itemize}
\tightlist
\item
  PSM 16 It is hard for me to get intensely interested in what is going
  on in my community. (Reversed)
\item
  PSM 23 I unselfishly contribute to my community.
\item
  PSM 30 Meaningful public service is very important to me.
\item
  PSM 34 I would prefer seeing public officials do what is best for the
  whole community even if it harmed my interests.
\item
  PSM 39 I consider public service my civic duty.
\end{itemize}
\end{minipage} \\
Affective motive & Compassion (aka ``patriotism of benevolence'') &
\begin{minipage}[t]{\linewidth}\raggedright
\begin{itemize}
\tightlist
\item
  PSM 2 I am rarely moved by the plight of the underprivileged.
  (Reversed)
\item
  PSM 3 Most social programs are too vital to do without.
\item
  PSM 4 It is difficult for me to contain my feelings when I see people
  in distress.
\item
  PSM 8 To me, patriotism includes seeing to the welfare of others.
\item
  PSM 10 I seldom think about the welfare of people whom I don't know
  personally. (Reversed)
\item
  PSM 13 I am often reminded by daily events about how dependent we are
  on one another.
\item
  PSM 24 I have little compassion for people in need who are unwilling
  to take the first step to help themselves. (Reversed)
\item
  PSM 40 There are few public programs that I wholeheartedly support.
  (Reversed)
\end{itemize}
\end{minipage} \\
Unclear & Self-sacrifice & \begin{minipage}[t]{\linewidth}\raggedright
\begin{itemize}
\item
  PSM 1 Making a difference in society means more to me than personal
  achievements.
\item
  PSM 5 I believe in putting duty before self.
\item
  PSM 6 Doing well financially is definitely more important to me than
  doing good deeds. (Reversed)
\item
  PSM 9 Much of what I do is for a cause bigger than myself.
\item
  PSM 12 Serving citizens would give me a good feeling even if no one
  paid me for it.
\item
  PSM 17 I feel people should give back to society more than they get
  from it.
\item
  PSM 19 I am one of those rare people who would risk personal loss to
  help someone else.
\item
  PSM 26 I am prepared to make enormous sacrifices for the good of
  society.
\end{itemize}
\end{minipage} \\
\end{longtable}

\newpage

\section{Appendix B}\label{appendix-b}

\subsection{Methodological Notes and Critiques Related to Literature on
Psychological
Empowerment}\label{methodological-notes-and-critiques-related-to-literature-on-psychological-empowerment}

\subsubsection{\texorpdfstring{On the Measurement of Psychological
Empowerment: Reivew Speicifc to Spreitzer
(\citeproc{ref-spreitzer1995}{1995})}{On the Measurement of Psychological Empowerment: Reivew Speicifc to Spreitzer (1995)}}\label{on-the-measurement-of-psychological-empowerment-reivew-speicifc-to-spreitzer1995}

As stated in the main text, the Psychological Empowerment scale
measurement instrument includes 12-items with the component dimensions
measured with three items each. In this case, psychological empowerment
is understood as a common second-order latent factor that underlies the
four latent component factors (e.g., meaning, impact, competence,
self-determination). Initial analysis of the measurement scale revealed
significant correlation between the component factors, suggesting a lack
of discriminant validity between the four dimensional components of
empowerment. Further research should be done to further refine the
theory and measurement of psychological empowerment
(\citeproc{ref-spreitzer1995}{Spreitzer 1995, 1458}).

Validity and reliability was assessed for the four components of
empowerment by examination of the overall Cronbach's alpha (\(\alpha\))
of the scale and the same was done for each component as well. The
measurement scales for the four components of empowerment were found to
be valid and moderately reliable. Second-order confirmatory factor
analysis (CFA) was used to assess the convergent and discriminant
validity of the empowerment measures in both samples. Second-order
confirmatory factor analysis is structurally equivalent to first-order
CFA: there is an unobservable latent factor (i.e., construct) theorized
to be influential upon observable referents, or indicators. For
second-order CFA, the latent factor is theorized to be a common higher
level factor influential upon a set of other lower level latent
component factors.

It's important to make a quick note on the survey sample and data
collection strategy used in this study to develop and validate the
measurement scale of psychological empowerment. Especially since this
scale is referenced and drawn upon often in subsequent research. Survey
data was collected for two distinct samples of a Fortune 50 industrial
company. The first sample (\(n = 393\)) comprised of managers from
across departments. Manager respondents were surveyed over a period of
three years until a 100\% response rate was achieved (this was not panel
data). This first primary sample consisted of mid-level managers,
\(93\%\) male, \(85\%\) White, average age of \(46\) years, an average
tenure of \(13\) years at the company, and average time in position of
\(3\) years. The extent that this sample is representative of people
with jobs, generally speaking, is probably pretty low. But more
importantly, the author was trying to develop and validate a novel
measure of empowerment by testing it on a sample that is likely to be
more empowered than the average worker. It is not clear why the sample
was confined to this group, but it seems like it would have made more
sense to randomly sample all employees.

In contrast, the second sample was surveyed at two different intervals
spanning five months in between survey administration, and consisted of
(non-managerial) employees (\(n = 128\)) of the same company. The author
reported a difference in measurement validity between the two samples;
the Cronbach's alpha for the overall measure of empowerment was
\(\alpha = 0.72\) for the primary industrial sample of managers, but
\(\alpha = 0.62\) for the second sample of employees. This calls into
question whether the equivalence assumption
(\citeproc{ref-pietryka2022}{Pietryka and MacIntosh 2022}) holds for
this measurement scale. Additionally, the primary sample of managers
were surveyed across a period of three years at the beginning of a
managerial development program. Seeing as how these managers all work in
the same company, it apparently wasn't considered that any of the
managers of the sample would talk to each other about the survey over
the period of three years, which potentially may have contaminated the
results of the study.

Using the new measure for psychological empowerment, Spreitzer
(\citeproc{ref-spreitzer1995}{1995}) tested whether the scale
significantly predicted two outcomes: managerial effectiveness and
innovative behavior. The former was defined as, ``the degree to which a
manager fulfills or exceeds work role expectations''
(\citeproc{ref-spreitzer1995}{1995, 1448}). Note that ``managerial
effectiveness'' does not only pertain to workers who are employed as
managers, but rather apparently refers to a kind of effectiveness which
would be described as ``managerial''. Spreitzer explains, ``Because, by
definition, empowered managers see themselves as competent and able to
influence their jobs and work environments in meaningful ways, they are
likely to proactively execute their job responsibilities by, for
instance, anticipating problems and acting independently, and hence are
likely to be seen as effective.'' (\citeproc{ref-spreitzer1995}{1995,
1448}). Innovative behavior was defined as ``the creation of something
new or different''.

It is unclear what survey items were used to measure these two
behavioral outcomes, but both managerial effectiveness and innovative
behavior was ascertained by questioning the subordinate employees of the
managers from the primary sample. The two behavioral outcomes were found
to be significantly related to empowerment, but the marginal fit of the
model was said to be likely due to multicollinearity between those two
behavioral outcomes.

Spreitzer also fit a Multiple Indicators Multiple Causes (MIMIC) model
to test whether certain factors were antecedent to empowerment. This
attempt to identify individual and contextual factors antecedent to
empowerment found partial support, but was not substantially convincing.
Many of the independent constructs theorized to influence empowerment
overlap in definition or description and are not independent in that
they potentially cause or influence each other. Each factor is
nevertheless theorized to have independent influence on empowerment. For
example, feedback (i.e., information about performance) and self-esteem
are considered as two independent factors that are antecedent to
empowerment. The former contextual, the latter a dispositional trait.
However, feedback on performance likely impacts self-esteem despite its
conception as a stable and enduring disposition. That raises issues of
endogeneity (two independent variables not so independent).

\subsubsection{\texorpdfstring{Psychological Empowerment as a Mediator:
Review Specific to Liden, Wayne, and Sparrowe
(\citeproc{ref-liden2000}{2000}) and Shapira-Lishchinsky and Tsemach
(\citeproc{ref-shapira-lishchinsky2014}{2014})}{Psychological Empowerment as a Mediator: Review Specific to Liden, Wayne, and Sparrowe (2000) and Shapira-Lishchinsky and Tsemach (2014)}}\label{psychological-empowerment-as-a-mediator-review-specific-to-liden2000-and-shapira-lishchinsky2014}

As stated in the main text, two of the three studies examining
psychological empowerment as a mediating variable appeared problematic
upon review. What follows are my initial impressions and critiques upon
review of Liden, Wayne, and Sparrowe (\citeproc{ref-liden2000}{2000})
and Shapira-Lishchinsky and Tsemach
(\citeproc{ref-shapira-lishchinsky2014}{2014}).

The earliest of these studies shows that the meaning and competence
dimensions of psychological empowerment have significant influence on
organizational commitment (\citeproc{ref-liden2000}{Liden, Wayne, and
Sparrowe 2000}). However, the cited study appears somewhat questionable
in terms of method and measurement of the relevant variables. The cited
study uses the term ``mediator variables'' to describe the empowerment
component dimensions, but the analysis conducted appears to have
consisted of multivariate OLS. The dependent variables were workplace
satisfaction, organizational commitment, and job performance, whereas
the independent variables included every component of empowerment,
leadership-member exchange (relationship with leadership) and
team-member exchange (relationship with coworkers), and job
characteristics.

The authors present three OLS models, with the last multivariate OLS
model being the most significant as it includes both the ``mediator
variables'' (empowerment components) and other independent variables all
together. Looking at this Table 4 (\citeproc{ref-liden2000}{Liden,
Wayne, and Sparrowe 2000, 413}), the only statistically significant
positively related independent variables to work satisfaction are job
characteristics and the meaning component of empowerment, whereas
competence shows a significant negative influence on work
satisfaction--on average and holding all else equal. This finding in
Table 4 suggests that an increase in feelings of competence in one's job
decreases work satisfaction. In addition, The measurement for work
satisfaction was measured by marking ``yes'' for adjectives connotating
satisfaction, but responses are coded so that a ``yes'' = 3, whereas a
``no'' = 0. This coding scheme would seemingly inflate the Cronbach's
alpha for the measure of work satisfaction and doesn't convey convincing
construct nor face validity. The measurement of the relevant variables
(save for the empowerment components) and the analysis of the study
makes the reported findings in Liden, Wayne, and Sparrowe
(\citeproc{ref-liden2000}{2000}) article appear dubious.

A relatively more recent study examined whether psychological
empowerment serves as a mediating variable between teachers perceptions
of authentic leadership and OCB and withdraw behaviors (e.g., lateness
and voluntary absence frequency)
(\citeproc{ref-shapira-lishchinsky2014}{Shapira-Lishchinsky and Tsemach
2014}). The essential findings of the research intuitively suggest that
when the boss is perceived as an `authentic leader', employees
(teachers) are increasingly motivated to perform above and beyond what
is required of them and report no intent to leave their current job.
Here, issues arise in the description and measurement of psychological
empowerment, which somewhat undermine the reported findings. The cited
authors employ a selective approach to the use of different measurement
scales, and when it came to psychological empowerment, chose to reduce
the component dimensions down to two and selected eight of the original
12 survey items.

The confirmatory factor analysis conducted by Shapira-Lishchinsky and
Tsemach (\citeproc{ref-shapira-lishchinsky2014}{2014}) yielded two
factors but the method is literally an analysis meant to confirm the
number of expected factors previously theorized. The authors
operationally conceived of psychological empowerment as a two
dimensional construct. The authors limited the number of factors down to
two, selected 8 out of 12 items from the original psychological
empowerment scale, and then conducted an analysis to confirm that survey
data sufficiently loads onto only two factors. The thing is, the
original psychological empowerment scale is comprised of four latent
factors derived from 12-items, where 3 items are used to measure each
component factor. The original conception of psychological empowerment
has four (latent) dimensions. Seeing as how the analysis by Spreitzer
(\citeproc{ref-spreitzer1995}{1995}) found that divergent validity was
lacking among the component factors, the selective use of various survey
items from the original scale is likely to only exacerbate such issues
of scale validity.

\newpage

\subsection*{References}\label{references}
\addcontentsline{toc}{subsection}{References}

\phantomsection\label{refs}
\begin{CSLReferences}{1}{1}
\bibitem[\citeproctext]{ref-ashforth1989}
Ashforth, Blake E. 1989. {``The Experience of Powerlessness in
Organizations.''} \emph{Organizational Behavior and Human Decision
Processes} 43(2): 207--42.
doi:\href{https://doi.org/10.1016/0749-5978(89)90051-4}{10.1016/0749-5978(89)90051-4}.

\bibitem[\citeproctext]{ref-bertelli2007}
Bertelli, Anthony M. 2007. {``Determinants of {Bureaucratic Turnover
Intention}: {Evidence} from the {Department} of the {Treasury}.''}
\emph{Journal of Public Administration Research and Theory: J-PART}
17(2): 235--58. \url{https://www.jstor.org/stable/4139743} (February 27,
2024).

\bibitem[\citeproctext]{ref-dahl1961}
Dahl, Robert A. 1961. \emph{Who {Governs}?: {Democracy} and {Power} in
an {American City}}. Yale University Press.

\bibitem[\citeproctext]{ref-dahl2015}
Dahl, Robert A., and Ian Shapiro. 2015. \emph{On Democracy}. Second
edition. New Haven: Yale University Press.
doi:\href{https://doi.org/10.12987/9780300233322}{10.12987/9780300233322}.

\bibitem[\citeproctext]{ref-deci2012}
Deci, Edward L., and Richard M. Ryan. 2012. {``Self-Determination
Theory.''} In \emph{Handbook of Theories of Social Psychology, {Vol}.
1}, Thousand Oaks, CA: Sage Publications Ltd, 416--36.
doi:\href{https://doi.org/10.4135/9781446249215.n21}{10.4135/9781446249215.n21}.

\bibitem[\citeproctext]{ref-gardner2000}
Gardner, Wendi L., Cynthia L. Pickett, and Marilynn B. Brewer. 2000.
{``Social Exclusion and Selective Memory: How the Need to Belong
Influences Memory for Social Events.''} \emph{Personality and Social
Psychology Bulletin} 26(4): 486--96.
doi:\href{https://doi.org/10.1177/0146167200266007}{10.1177/0146167200266007}.

\bibitem[\citeproctext]{ref-hackman1974}
Hackman, J. Richard, and Greg R. Oldham. 1974. \emph{The {Job Diagnostic
Survey}: {An Instrument} for the {Diagnosis} of {Jobs} and the
{Evaluation} of {Job Redesign Projects}}. Yale University. Technical
Report. \url{https://eric.ed.gov/?id=ED099580} (March 28, 2024).

\bibitem[\citeproctext]{ref-hackman1975}
Hackman, J. Richard, and Greg R. Oldham. 1975. {``Development of the
{Job Diagnostic Survey}.''} \emph{Journal of Applied Psychology} 60(2):
159--70. doi:\href{https://doi.org/10.1037/h0076546}{10.1037/h0076546}.

\bibitem[\citeproctext]{ref-higgins2021}
Higgins, E Tory, and Emily Nakkawita. 2021. {``Motives {Working
Together}: {Value}, {Truth}, and {Control} in {Goal Selection} and
{Pursuit}.''} In \emph{Social Psychology : Handbook of Basic
Principles}, eds. Paul A. M. van Lange, Arie W. Kruglanski, and E Tory
Higgins. New York: The Guilford Press, 85--105.
\url{https://search.ebscohost.com/login.aspx?direct=true&scope=site&db=nlebk&db=nlabk&AN=2645023}.

\bibitem[\citeproctext]{ref-hull1943}
Hull, C. L. 1943. \emph{Principles of Behavior: An Introduction to
Behavior Theory}. Oxford, England: Appleton-Century.

\bibitem[\citeproctext]{ref-kruglanski2017}
Kruglanski, Arie W., Ayelet Fishbach, Kaitlin Woolley, Jocelyn Belanger,
Marina Chernikova, Erica Molinario, and Antonio Pierro. 2017. {``A
{Structural Model} of {Intrinsic Motivation}: {On} the {Psychology} of
{Means-Ends Fusion}.''}
doi:\href{https://doi.org/10.31234/osf.io/kdct8}{10.31234/osf.io/kdct8}.

\bibitem[\citeproctext]{ref-kruglanski2022}
Kruglanski, Arie W., Erica Molinario, Katarzyna Jasko, David Webber, N.
Pontus Leander, and Antonio Pierro. 2022. {``Significance-{Quest
Theory}.''} \emph{Perspectives on Psychological Science} 17(4):
1050--71.
doi:\href{https://doi.org/10.1177/17456916211034825}{10.1177/17456916211034825}.

\bibitem[\citeproctext]{ref-liden2000}
Liden, Robert C., Sandy J. Wayne, and Raymond T. Sparrowe. 2000. {``An
Examination of the Mediating Role of Psychological Empowerment on the
Relations Between the Job, Interpersonal Relationships, and Work
Outcomes.''} \emph{Journal of Applied Psychology} 85(3): 407--16.
doi:\href{https://doi.org/10.1037/0021-9010.85.3.407}{10.1037/0021-9010.85.3.407}.

\bibitem[\citeproctext]{ref-mcleod2022}
Mcleod, Saul. 2022. {``Maslow's Hierarchy of Needs.''}
\url{https://www.simplypsychology.org/maslow.html}.

\bibitem[\citeproctext]{ref-perry1996}
Perry, James L. 1996. {``Measuring {Public Service Motivation}: {An
Assessment} of {Construct Reliability} and {Validity}.''} \emph{Journal
of Public Administration Research and Theory: J-PART} 6(1): 5--22.
\url{https://www.jstor.org/stable/1181620} (February 27, 2024).

\bibitem[\citeproctext]{ref-pietryka2022}
Pietryka, Matthew T., and Randall C. MacIntosh. 2022. {``{ANES Scales
Often Do Not Measure What You Think They Measure}.''} \emph{The Journal
of Politics} 84(2): 1074--90.
doi:\href{https://doi.org/10.1086/715251}{10.1086/715251}.

\bibitem[\citeproctext]{ref-pitts2011}
Pitts, David, John Marvel, and Sergio Fernandez. 2011. {``So {Hard} to
{Say Goodbye}? {Turnover Intention} Among {U}.{S}. {Federal
Employees}.''} \emph{Public Administration Review} 71(5): 751--60.
\url{https://www.jstor.org/stable/23017442} (February 27, 2024).

\bibitem[\citeproctext]{ref-ryan2000}
Ryan, Richard M., and Edward L. Deci. 2000. {``Intrinsic and {Extrinsic
Motivations}: {Classic Definitions} and {New Directions}.''}
\emph{Contemporary Educational Psychology} 25(1): 54--67.
doi:\href{https://doi.org/10.1006/ceps.1999.1020}{10.1006/ceps.1999.1020}.

\bibitem[\citeproctext]{ref-ryan2017}
Ryan, Richard M., and Edward L. Deci. 2017. \emph{Self-Determination
Theory : Basic Psychological Needs in Motivation, Development, and
Wellness}. New York: Guilford Press.
\url{http://books.google.com/books?vid=isbn9781462528769}.

\bibitem[\citeproctext]{ref-ryan2021}
Ryan, Richard M., and Edward L. Deci. 2021. {``Human {Autonomy} in
{Social Psychology}: {A Self-Determination Theory Perspective}.''} In
\emph{Social Psychology : Handbook of Basic Principles}, eds. Paul A. M.
van Lange and Arie W. Kruglanski. New York: The Guilford Press.

\bibitem[\citeproctext]{ref-shapira-lishchinsky2014}
Shapira-Lishchinsky, Orly, and Sigalit Tsemach. 2014. {``Psychological
{Empowerment} as a {Mediator Between Teachers}' {Perceptions} of
{Authentic Leadership} and {Their Withdrawal} and {Citizenship
Behaviors}.''} \emph{Educational Administration Quarterly} 50(4):
675--712.
doi:\href{https://doi.org/10.1177/0013161X13513898}{10.1177/0013161X13513898}.

\bibitem[\citeproctext]{ref-skinner1965}
Skinner, Burrhus Frederic. 1965. \emph{Science and Human Behavior}.
{Simon and Schuster}.

\bibitem[\citeproctext]{ref-spreitzer1995}
Spreitzer, Gretchen M. 1995. {``Psychological {Empowerment} in the
{Workplace}: {Dimensions}, {Measurement}, and {Validation}.''} \emph{The
Academy of Management Journal} 38(5): 1442--65.
doi:\href{https://doi.org/10.2307/256865}{10.2307/256865}.

\bibitem[\citeproctext]{ref-taylor2013}
Taylor, Jeannette. 2013. {``Goal {Setting} in the {Australian Public
Service}: {Effects} on {Psychological Empowerment} and {Organizational
Citizenship Behavior}.''} \emph{Public Administration Review} 73(3):
453--64. \url{https://www.jstor.org/stable/42002950} (February 27,
2024).

\bibitem[\citeproctext]{ref-tett1993}
Tett, Robert P., and John P. Meyer. 1993. {``{JOB SATISFACTION},
{ORGANIZATIONAL COMMITMENT}, {TURNOVER INTENTION}, {AND TURNOVER}: {PATH
ANALYSES BASED ON META-ANALYTIC FINDINGS}.''} \emph{Personnel
Psychology} 46(2): 259--93.
\url{https://search.ebscohost.com/login.aspx?direct=true&db=eue&AN=9401104761&site=ehost-live}.

\bibitem[\citeproctext]{ref-thomas1990}
Thomas, Kenneth W., and Betty A. Velthouse. 1990. {``Cognitive
{Elements} of {Empowerment}: {An} "{Interpretive}" {Model} of {Intrinsic
Task Motivation}.''} \emph{The Academy of Management Review} 15(4):
666--81. doi:\href{https://doi.org/10.2307/258687}{10.2307/258687}.

\end{CSLReferences}



\end{document}
